\documentclass[10pt,a4paper]{amsart}
\usepackage[margin=19mm]{geometry}
\usepackage{amsmath,amssymb,mathtools}
\usepackage[scr=rsfs,cal=euler]{mathalfa}
\usepackage{xfrac}
\usepackage[unicode,hidelinks,bookmarks=false]{hyperref}
\newcommand{\ie}{i.e.\ }
\newcommand{\cf}{cf.\ }
\newcommand{\resp}{respectively\ }
\newcommand{\Subsection}{\subsection}
\providecommand{\nobreakdash}{\nobreak\hbox{-}}
\setlength{\emergencystretch}{2em}
\multlinegap0pt
\usepackage{amssymb}
\newtheorem{lemma}{Lemma}
\newcommand{\Nzero}{\mathbb N_0}
\title{\Large On a problem of Nathanson related to minimal asymptotic bases and maximal asymptotic nonbases}
\author{Shi-Qiang Chen}
\date{Source: arXiv:2609.07361v1 (CC BY 4.0)}
\begin{document}
\maketitle

\noindent\emph{Source and licence.} \Large On a problem of Nathanson related to minimal asymptotic bases and maximal asymptotic nonbases. Version arXiv:2609.07361v1 (CC BY 4.0), distributed by arXiv under the Creative Commons Attribution 4.0 International licence, http://creativecommons.org/licenses/by/4.0/, as recorded in the arXiv licence metadata for this exact version and shown on the abstract page https://arxiv.org/abs/2609.07361v1. Full licence notice: \texttt{licenses/arxiv-2609.07361-CC-BY-4.0.txt}.\par\smallskip

\noindent\textit{Editorial scope.} Counted unit: the article's lemma lem3 (source0.tex lines 324--353). The article's complete original proof of it is delivered with it (source0.tex lines 355--577, one \texttt{proof} environment of 223 lines). No mathematics is written, repaired or re-derived: the statement and the proof are the article's own lines, in the article's own notation, and no retained line is substituted or re-flowed.  Retained uncounted as the source-local context the proof needs: lines 79--89.  Nothing was omitted from any retained span, so the package carries no omission record.  No retained line contains a citation, so the wrapper carries no bibliography and no bibliography entry.  No retained line contains a figure environment, an image-inclusion command or drawing code, so no image is delivered and the package declares no asset dependency.  Editorial formatting only, in addition: an AMS \texttt{amsart} preamble that restates the article's own declarations and macros used by the retained lines, namely the environments \texttt{lemma} on separate counters, together with the standard packages those lines need.  The retained environments print in this document as Lemma 1 in order of appearance, whereas the article prints the same items with the numbering of its own full document.  The complete native source of the article as distributed in the arXiv e-print for this exact version is bundled unchanged as \texttt{sources/209621/source0.tex}.
\begin{lemma}\label{lem1}
Let $S\subseteq\Nzero$, let $p,u,v$ be nonnegative integers with $p<u\le v$. If $p\in S$, $[u,v]\subseteq S$ and $v\ge 2u-p-1$,
then for every integer $m\ge2$, we have
\[
 [u+(m-1)p,mv]\subseteq mS.
\]
In particular, if $0\in S$ and $v\ge2u-1$, then
\[
 [u,mv]\subseteq mS.
\]
\end{lemma}

\begin{lemma}\label{lem3}
Let $T$ and $U$ be two positive integers with $T\ge T_0$ and $T=2U-2$, and let $C$ and $D$ be two sets such that
\[
 0\in D,\qquad 1\in C,\qquad [d,d+2h]\subseteq D,\qquad C\sqcup D=[0,T],\qquad [U,T]\subseteq C.
\]
Put
\[
 L=2T,\qquad M=(2h-3)T+1,\qquad W=2hM+1.
\]
Fix $c\in C$ and put
\[
 y_c=W-(h-2)L-c.
\]
Put
\[
 H=\{L\}\cup[M,2M]\cup\{y_c\}\cup[W-h+2,W].
\]
We define
\[
 E=C\cup H,\qquad F=D\cup([T+1,W]\setminus H).
\]
Then
\begin{equation}\label{eq5}
[T+1,W]\subseteq hE\cap hF,
\end{equation}
and
\begin{equation}\label{eq6}
W\notin h(E\setminus\{c\}).
\end{equation}
\end{lemma}

\begin{proof}
Since $0\in D$, $c\in C$ and $C\sqcup D=[0,T]$, it follows that
\[
 1\le c\le T.
\]
Noting that
\[
 (h-2)L+T=(2h-3)T=M-1,
\]
we have
\[
 y_c\ge W-(M-1)=(2h-1)M+2>2M.
\]
Moreover,
\[
 W-h+2-y_c=(h-2)L+c-h+2>0.
\]
Consequently,
\[
 2M<y_c<W-h+2,
\]
so the four pieces occurring in the definition of $H$ are pairwise disjoint.

By the definition of $y_c$, we have
\[
 W=y_c+\underbrace{L+\cdots+L}_{h-2\text{ copies}}+c\in hE.
\]
We first prove \eqref{eq6}. Noting that $0\in D$ and $C\sqcup D=[0,T]$, we have $0\notin E$. It follows that $e>0$ for every $e\in E$. Write
\[W=w_1+w_2+\ldots+w_h,~~~ w_i\in E,~~~i=1,2,\ldots,h.
\]
If there exists an integer $i_0\in\{1,2,\ldots,h\}$ such that
\[
 w_{i_0}\in[W-h+2,W],
\]
then 
\[
 W\ge w_{i_0}+h-1\ge W+1,
\]
a contradiction. Thus, $w_i\not\in [W-h+2,W]$ for $i=1,2,\ldots,h$.

Since
\[
 y_c-\frac W2\ge (h-1)M+\frac32>0,
\]
it follows that
\[
\sum\limits_{i=1}^{h}|\{w_i\}\cap\{y_c\}|\leq 1.
\] 
If \[
\sum\limits_{i=1}^{h}|\{w_i\}\cap\{y_c\}|=0,
\] 
then 
\[
 W\le 2hM=W-1,
\]
a contradiction. Therefore,
\[
\sum\limits_{i=1}^{h}|\{w_i\}\cap\{y_c\}|=1.
\] 
It follows that $W-y_c\in (h-1)E$ and
\[
 W-y_c=(h-2)L+c\le(h-2)L+T=M-1.
\]
By the definition of $E$, we have
\[
 E\cap[0,M-1]=C\cup\{L\},
\]
and so \[
W-y_c\in (h-1)(C\cup\{L\}).
\] 
Let $r$ be the number of copies of $L$ among the remaining $h-1$ summands. If $r\le h-3$, then
\[
W-y_c\leq rL+(h-1-r)T=(h+r-1)T\le(2h-4)T=(h-2)L,
\]
which is impossible because $c\ge1$. If $r=h-1$, then
\[
 (h-2)L+c=(h-1)L,
\]
which is impossible because $c\le T<L$. Thus $r=h-2$, so there exists an integer $i_1\in\{1,2,\ldots,h\}$ such that $w_{i_1}=c$, which proves \eqref{eq6}.

We next prove \eqref{eq5}. Noting that $1\in E$, $[U,T]\subseteq E$ and 
\[
 T=2U-2=2U-1-1,
\]
Lemma~\ref{lem1} applies with $S=E$, $p=1$, $u=U$, $v=T$ and $m=h-k$ for every $0\le k\le h-2$.
Then
\[
 I_k=[2kT+U+h-k-1,(h+k)T]\subseteq hE.
\]
Noting that $U=(T+2)/2$ and $T\ge T_0$, we have $U\ge h+2$, and hence
\[
 \min I_0=U+h-1\le T+1.
\]
For $0\le k\le h-3$, by $T=2U-2$ and $U\ge h+2$, we have
\[
 U+h-k-3\le(h-k-2)T.
\]
Moreover,
\[
 \max I_{h-2}=(2h-2)T\ge(2h-3)T=M-1.
\]
Consequently,
\[
 [T+1,M-1]\subseteq hE.
\]
Noting that $T,1,L\in E$, we have
\[
 M=\underbrace{L+\cdots+L}_{h-2\text{ copies}}+T+1\in hE.
\]
For $1\le j\le h-2$, noting that
\[
 T+1+j\in[2U,2T]=2[U,T],
\]
we have
\[
 M+j=(h-2)L+(T+1+j)\in hE.
\]
Finally, Lemma~\ref{lem1} applied to $S=E$, $p=1$, $u=M$ and $v=2M$ gives
\[
 [M+h-1,2hM]=[M+h-1,W-1]\subseteq hE.
\]
It follows from $W\in hE$ that
\[
 [T+1,W]\subseteq hE.
\]
Noting that $e\in F$ for every $e\in [T+1,W]\setminus H$ and $0\in F$, we have $e\in hF$ for every $e\in [T+1,W]\setminus H$. Thus it is enough to cover the four pieces of $H$.

Noting that $d\in D$, $d<T$ and $L=2T$, we have
\[
 T<L-d<L.
\]
Hence, $L-d\in[T+1,L-1]\subseteq F$, and so
\[
 L=(L-d)+d+(h-2)\cdot0\in hF.
\]

Next, put
\[
 J=[L+1,M-1]=[2T+1,M-1]\subseteq F.
\]
Suppose $h\ge4$. Since
\[
 4T+2\le M,
\]
it follows that
\[
 [M,2M-2]\subseteq2J=[4T+2,2M-2]\subseteq hF.
\]
Moreover,
\[
 6T+3\le2M-1\le2M\le3M-3,
\]
so
\[
 2M-1,2M\in3J\subseteq hF.
\]
Thus,
\[
 [M,2M]\subseteq hF.
\]


If $h=3$, then $M=3T+1$ and $J=[2T+1,3T]$. It follows that
\[
 [T+1,2T-1]+J=[3T+2,5T-1]\subseteq hF
\]
and
\[
 2J=[4T+2,6T]\subseteq hF.
\]
Thus, $[M+1,2M-2]\subseteq hF$. Since 
\[
 d,d+1,d+2\in[d,d+2h]\subseteq D,
\]
while $M-d,3T-d\in J$, it follows that
\[
 M=d+(M-d)+(h-2)\cdot0\in hF,
\]
and
\[
 2M-1=3T+(3T-d)+(d+1)\in hF,
\]
\[
 2M=3T+(3T-d)+(d+2)\in hF.
\]
Hence, $[M,2M]\subseteq hF$.

Since $d+2h=12h$ and $T\ge T_0$, it follows that
\[
 y_c-(2M+d+2h)\ge(2h-3)M+2-12h>0.
\]
Thus for every $b\in[d,d+2h]$, we have
\[
 2M<y_c-b<y_c<W-h+2.
\]
By the construction, $y_c-b\in F$, we have
\[
 y_c=(y_c-b)+b+(h-2)\cdot0\in hF.
\]
Let $z\in[W-h+2,W]$. Choose $b\in[d,d+2h]$ such that $z-b\ne y_c$. Since $d=10h>h-2$, it follows that
\[
 z-b\le W-d<W-h+2.
\]
On the other hand,
\[
 z-b\ge W-h+2-(d+2h)=2hM+3-13h>2M.
\]
Therefore,
\[
 2M<z-b<W-h+2,
\]
and, by the choice of $b$, the integer $z-b$ does not belong to any of the four pieces of $H$. Hence, $z-b\in F$ and
\[
 z=(z-b)+b+(h-2)\cdot0\in hF.
\]
It follows that
\[
 [T+1,W]\subseteq hF.
\]
Combining the two inclusions proves \eqref{eq5}. 

This completes the proof of Lemma \ref{lem3}.
\end{proof}

\end{document}
